\documentclass[11pt]{article}

\usepackage[a4paper,margin=24mm]{geometry}
\usepackage{amsmath,amssymb,mathtools}
\usepackage{booktabs}
\usepackage{float}
\usepackage{graphicx}
\usepackage{microtype}
\usepackage{tabularx}
\usepackage{placeins}
\usepackage[hidelinks]{hyperref}
\usepackage{xcolor}

\newcommand{\trans}{\mathsf{T}}
\newcommand{\shortp}{\mathrm{s}}
\newcommand{\longp}{\ell}
\newcommand{\Frob}{\mathrm{F}}
\newcommand{\SO}{\mathrm{SO}(2)}
\newcommand{\SLZ}{\mathrm{SL}(2,\mathbb{Z})}
\newcommand{\Uminus}{U_{-}}
\newcommand{\Lminus}{L_{-}}
\newcommand{\Lplus}{L_{+}}
\newcommand{\Uplus}{U_{+}}

\definecolor{shortblue}{HTML}{2171B5}
\definecolor{longorange}{HTML}{E67E22}

\title{Two Elastic-Reduction Quadrants and the\
Plastic--Elastic Decomposition of a Simple Shear}
\author{Technical note}
\date{14 July 2026}

\begin{document}
\maketitle

\begin{abstract}
For the simple shear
\(
 \mathbf F=\left(\begin{smallmatrix}1&1.6\\0&1\end{smallmatrix}\right)
\), the unit-shear breadth-first search reaches two points in the elastic
domain.  The shortest and longer routes produce visibly different matrices
\(\mathbf M\), and hence apparently different multiplicative decompositions
\(\mathbf F=\mathbf F_E\mathbf F_P\).  The main observation of this note is
that the two decompositions differ only by a $90^\circ$ rotation of the
intermediate basis.  They have the same spatial elastic metric
\(\mathbf B_E\), the same reference plastic metric \(\mathbf C_P\), the same
elastic singular values, energy, and Cauchy stress.  Thus the two elastic
quadrants are physically equivalent unless an additional history or
orientation variable fixes the intermediate frame.  A direct comparison in
the original coordinate frame does provide such a preference: the short-path
plastic factor retains the positive horizontal shear of the total
deformation, while the long-path factor looks like a vertical shear.  The
short factor is also much closer to the original \(\mathbf F\).  This is a
meaningful frame-anchoring or continuity criterion, although the preference
disappears if the intermediate frame is allowed to rotate by its admissible
$90^\circ$ lattice symmetry.
\end{abstract}

\section{The two multiplicative decompositions}

The accumulated integer matrix from elastic reduction acts on the right,
\begin{equation}
 \mathbf C_E=\mathbf M^\trans\mathbf C\mathbf M,
 \qquad
 \mathbf F_E=\mathbf F\mathbf M.
\end{equation}
Using the convention in the decomposition code,
\begin{equation}
 \boxed{
 \mathbf F=\mathbf F_E\mathbf F_P,
 \qquad
 \mathbf F_P=\mathbf M^{-1}.}
 \label{eq:decomposition}
\end{equation}
For the present \(\mathbf F\), the shortest route has two unit shears,
whereas the first route to the other elastic quadrant has three:
\begin{align}
 \text{short: }&\Uminus\Uminus,
 &\mathbf M_{\shortp}&=
 \begin{pmatrix}1&-2\\0&1\end{pmatrix},
 \\
 \text{long: }&\Uminus\Lminus\Uplus,
 &\mathbf M_{\longp}&=
 \begin{pmatrix}2&1\\-1&0\end{pmatrix}.
\end{align}
Both factorizations reproduce the total deformation exactly.  Their elastic
and plastic factors are
\begin{align}
 \mathbf F_E^{(\shortp)}&=
 \begin{pmatrix}1&-0.4\\0&1\end{pmatrix},
 &
 \mathbf F_P^{(\shortp)}&=
 \begin{pmatrix}1&2\\0&1\end{pmatrix},
 \\
 \mathbf F_E^{(\longp)}&=
 \begin{pmatrix}0.4&1\\-1&0\end{pmatrix},
 &
 \mathbf F_P^{(\longp)}&=
 \begin{pmatrix}0&-1\\1&2\end{pmatrix}.
\end{align}
Figure~\ref{fig:factors} puts these factors beside the matrices to which they
should naturally be compared: \(\mathbf F_P\) beside the original
\(\mathbf F\), and \(\mathbf F_E\) beside the undeformed benchmark
\(\mathbf I\).  To distinguish direct componentwise comparison from comparison
after an admissible change of intermediate basis, define
\begin{equation}
 \mathcal G=\{\mathbf R(k\pi/2):k=0,1,2,3\}
 =\SO\cap\SLZ.
 \label{eq:latticegroup}
\end{equation}
The corresponding numerical tests are collected in
Table~\ref{tab:decompositions}.

\begin{figure}[htbp]
 \centering
 \includegraphics[width=0.98\textwidth]{figures/elastic_plastic_factors.pdf}
 \caption{The original deformation and its two decompositions, together with
 the corresponding unit-shear BFS paths in the Poincar\'e disk.  In every
 matrix panel the solid vector is the first column \(\mathbf A_{:1}\), and the
 dashed vector is the second column \(\mathbf A_{:2}\).  The original
 \(\mathbf F\) and the short-path \(\mathbf F_P\) are both positive horizontal
 shears, whereas the long-path \(\mathbf F_P\) appears vertically oriented.
 The differently oriented elastic factors compensate exactly, so both
 products \(\mathbf F_E\mathbf F_P\) recover the same \(\mathbf F\).  Rotating
 the long-path intermediate basis back by \(90^\circ\) makes its plastic
 factor coincide with the short-path factor.}
 \label{fig:factors}
\end{figure}

\begin{table}[htbp]
\centering
\caption{Comparison with the original total deformation and the undeformed
benchmark.  ``Fixed'' means that all component matrices are read in the
displayed axes; ``aligned'' allows the lattice rotations in
\(\mathcal G\).}
\label{tab:decompositions}
\footnotesize
\renewcommand{\arraystretch}{1.28}
\begin{tabularx}{\textwidth}{@{}>{\raggedright\arraybackslash}p{31mm} >{\centering\arraybackslash}p{25mm} >{\centering\arraybackslash}p{25mm} >{\centering\arraybackslash}p{25mm} X@{}}
\toprule
Quantity & Original or benchmark & \textcolor{shortblue}{Shortest route} &
\textcolor{longorange}{Longer route} & Interpretation \\
\midrule
Plastic matrix &
$\mathbf F=\left(\begin{smallmatrix}1&1.6\\0&1\end{smallmatrix}\right)$ &
$\left(\begin{smallmatrix}1&2\\0&1\end{smallmatrix}\right)$ &
$\left(\begin{smallmatrix}0&-1\\1&2\end{smallmatrix}\right)$ &
Only the short factor retains the displayed upper-shear form \\
Fixed \(\lVert\mathbf F_P-\mathbf F\rVert_{\Frob}\) & $0$ & $0.400$ & $3.124$ &
Strongly favors the short route \\
Aligned \(\min_{\mathbf R\in\mathcal G}
\lVert\mathbf R\mathbf F_P-\mathbf F\rVert_{\Frob}\) & $0$ & $0.400$ & $0.400$ &
Tie after intermediate-frame alignment \\
Reference-metric \(\lVert\mathbf C_P-\mathbf C\rVert_{\Frob}\) & $0$ & $1.547$ & $1.547$ &
Objective comparison with the original \(\mathbf C\) is tied \\
Elastic matrix & $\mathbf I$ &
$\left(\begin{smallmatrix}1&-0.4\\0&1\end{smallmatrix}\right)$ &
$\left(\begin{smallmatrix}0.4&1\\-1&0\end{smallmatrix}\right)$ &
Only the short factor remains near the displayed axes \\
Fixed \(\lVert\mathbf F_E-\mathbf I\rVert_{\Frob}\) & $0$ & $0.400$ & $1.833$ &
Favors the short route \\
Objective \(\lVert\mathbf B_E-\mathbf I\rVert_{\Frob}\) & $0$ & $0.588$ & $0.588$ &
Spatial elastic comparison is tied \\
BFS depth & --- & $2$ & $3$ & Favors the short route if each move has positive cost \\
\bottomrule
\end{tabularx}
\end{table}
\FloatBarrier

\subsection{Original-frame preference and its limit}

With the intermediate axes identified with the original coordinate axes, the
preference is unambiguous.  The total deformation and the short-path plastic
factor are positive upper-triangular shears of amplitudes \(1.6\) and \(2\),
respectively.  Their difference is only the residual shear \(0.4\).  The
long-path plastic factor has a different displayed orientation and is much
farther from the total deformation.  The elastic comparison tells the same
story: the short-path factor is close to \(\mathbf I\), whereas the long-path
factor contains an additional quarter turn.

The limitation is that these direct matrix differences silently identify the
intermediate vector space with the original reference and spatial coordinate
spaces.  The two integer matrices satisfy
\begin{equation}
 \mathbf M_{\longp}=\mathbf M_{\shortp}\mathbf Q,
 \qquad
 \mathbf Q=
 \begin{pmatrix}0&1\\-1&0\end{pmatrix}
 =\mathbf R(-\pi/2),
 \qquad
 \mathbf Q^\trans\mathbf Q=\mathbf I.
 \label{eq:Qrelation}
\end{equation}
Consequently,
\begin{equation}
 \mathbf F_E^{(\longp)}=\mathbf F_E^{(\shortp)}\mathbf Q,
 \qquad
 \mathbf F_P^{(\longp)}=\mathbf Q^\trans\mathbf F_P^{(\shortp)}.
 \label{eq:gaugepair}
\end{equation}
The factors therefore change, but the rotation cancels from their product:
\begin{equation}
 (\mathbf F_E^{(\shortp)}\mathbf Q)
 (\mathbf Q^\trans\mathbf F_P^{(\shortp)})
 =\mathbf F_E^{(\shortp)}\mathbf F_P^{(\shortp)}=\mathbf F.
\end{equation}
The apparently vertical long-path plastic shear is therefore exactly the
short-path horizontal shear written in a rotated intermediate frame:
\begin{equation}
 \mathbf Q\mathbf F_P^{(\longp)}
 =\mathbf F_P^{(\shortp)}.
 \label{eq:alignplastic}
\end{equation}
Accordingly, the fixed-frame plastic discrepancy selects the short path, but
the discrepancy after lattice-compatible alignment does not:
\begin{align}
 \lVert\mathbf F_P^{(\shortp)}-\mathbf F\rVert_{\Frob}
 &=0.400,
 &
 \lVert\mathbf F_P^{(\longp)}-\mathbf F\rVert_{\Frob}
 &=3.124,
 \\
 \min_{\mathbf R\in\mathcal G}
 \lVert\mathbf R\mathbf F_P^{(\shortp)}-\mathbf F\rVert_{\Frob}
 &=
 \min_{\mathbf R\in\mathcal G}
 \lVert\mathbf R\mathbf F_P^{(\longp)}-\mathbf F\rVert_{\Frob}
 =0.400.
 \label{eq:alignedplasticdistance}
\end{align}
The minimizing rotations are \(\mathbf I\) for the short path and
\(\mathbf Q\) for the long path.  The same conclusion follows from a
comparison that never identifies the intermediate axes: since both routes
have the same \(\mathbf C_P\),
\begin{equation}
 \lVert\mathbf C_P^{(\shortp)}-\mathbf C\rVert_{\Frob}
 =
 \lVert\mathbf C_P^{(\longp)}-\mathbf C\rVert_{\Frob}
 =1.547.
 \label{eq:plasticmetricdistance}
\end{equation}
This is the standard rotational freedom of the intermediate configuration:
for any \(\mathbf Q\in\SO\),
\(
 \mathbf F_E\mathbf F_P
 =(\mathbf F_E\mathbf Q)(\mathbf Q^\trans\mathbf F_P)
\).
Here the admissible rotation is additionally lattice-compatible,
\(\mathbf Q\in\SO\cap\SLZ\).

The quantities insensitive to this intermediate rotation are identical:
\begin{align}
 \mathbf B_E^{(\shortp)}
 &=\mathbf F_E^{(\shortp)}\mathbf F_E^{(\shortp)\trans}
 =\mathbf B_E^{(\longp)}
 =\begin{pmatrix}1.16&-0.4\\-0.4&1\end{pmatrix},
 \\
 \mathbf C_P^{(\shortp)}
 &=\mathbf F_P^{(\shortp)\trans}\mathbf F_P^{(\shortp)}
 =\mathbf C_P^{(\longp)}
 =\begin{pmatrix}1&2\\2&5\end{pmatrix}.
 \label{eq:equalmetrics}
\end{align}
The two \(\mathbf F_E\) have the same singular values,
\(1.219804\) and \(0.819804\); the two \(\mathbf F_P\) have singular values
\(2.414214\) and \(0.414214\).  They also give the same energy,
\(0.0864665\), and the same Cauchy stress,
\begin{equation}
 \boldsymbol\sigma=
 \begin{pmatrix}
 -0.0322271&-0.0809712\\
 -0.0809712&0.0322271
 \end{pmatrix}.
\end{equation}
This explains why Cauchy stress cannot select between the quadrants: the
spatial elastic metric \(\mathbf B_E\), which carries the objective elastic
stretch, is exactly the same.

\section{What changes for even longer BFS paths?}

At maximum BFS depth five, seven terminal paths are found.  For this example,
every accumulated matrix belongs to one of only four classes:
\begin{equation}
 \mathbf M\in
 \{\mathbf M_{\shortp},-\mathbf M_{\shortp},
   \mathbf M_{\longp},-\mathbf M_{\longp}\}.
 \label{eq:Mclasses}
\end{equation}
The observed paths are summarized in Table~\ref{tab:paths}.

\begin{table}[htbp]
\centering
\caption{All terminal BFS paths through depth five.}
\label{tab:paths}
\small
\renewcommand{\arraystretch}{1.18}
\begin{tabularx}{0.88\textwidth}{@{}c l c X@{}}
\toprule
Quadrant & Unit-shear word & Depth & Accumulated matrix class \\
\midrule
1 & $\Uminus\,\Uminus$ & 2 & $\mathbf M_{\shortp}$ \\
1 & $\Lminus\,\Uplus\,\Uplus\,\Lminus$ & 4 & $-\mathbf M_{\shortp}$ \\
2 & $\Uminus\,\Lminus\,\Uplus$ & 3 & $\mathbf M_{\longp}$ \\
2 & $\Lminus\,\Uplus\,\Lplus$ & 3 & $\mathbf M_{\longp}$ \\
2 & $\Uplus\,\Lminus\,\Uplus\,\Lplus\,\Lplus$ & 5 & $\mathbf M_{\longp}$ \\
2 & $\Uminus\,\Lplus\,\Uminus\,\Lplus\,\Lplus$ & 5 & $-\mathbf M_{\longp}$ \\
2 & $\Lplus\,\Uminus\,\Lplus\,\Lplus\,\Lplus$ & 5 & $-\mathbf M_{\longp}$ \\
\bottomrule
\end{tabularx}
\end{table}

Thus the longer paths do not create new plastic or elastic stretch states in
this case.  Replacing \(\mathbf M\) by \(-\mathbf M\) gives
\(\mathbf F_E\mapsto-\mathbf F_E\) and
\(\mathbf F_P\mapsto-\mathbf F_P\).  Their product, \(\mathbf C_E\),
\(\mathbf B_E\), and \(\mathbf C_P\) are unchanged.  It is a $180^\circ$
rotation of the oriented intermediate basis.  Moving between the two elastic
quadrants adds the $90^\circ$ rotation in Equation~\eqref{eq:Qrelation}.

This four-class conclusion is specific to the present generic metric and the
paths searched here.  A metric with additional lattice symmetries can have a
larger stabilizer.  The robust statement is that two factorizations are
physically equivalent whenever their accumulated matrices differ by a
consistently applied orthogonal change of intermediate basis.

If the intermediate frame is anchored to the original axes, the four matrix
classes are no longer equally convenient.  Write
\(d_P=\lVert\mathbf F_P-\mathbf F\rVert_{\Frob}\) and
\(d_E=\lVert\mathbf F_E-\mathbf I\rVert_{\Frob}\).  Their scores are
\begin{equation}
 \begin{array}{c|cccc}
 \mathbf M & \mathbf M_{\shortp} & -\mathbf M_{\shortp}
 & \mathbf M_{\longp} & -\mathbf M_{\longp} \\ \hline
 (d_P,d_E) & (0.400,0.400) & (4.578,2.857)
 & (3.124,1.833) & (3.370,2.227)
 \end{array}
 \label{eq:fourclassscores}
\end{equation}
Thus the anchored-frame criterion uniquely selects
\(\mathbf M_{\shortp}\), including over the still longer paths.  If rotations
in \(\mathcal G\) are first aligned away, all four classes again represent the
same stretch state.

\section{Basis-rotation context}

``Basis transformation'' is a reasonable name for the operation in the
question.  More precisely, it is a \emph{passive orthogonal change of the
reference basis}.  If the new basis is represented by
\(\mathbf Q(\theta)=\mathbf R(\theta)\), the components transform as
\begin{equation}
 \mathbf F(\theta)=\mathbf F(0)\mathbf Q(\theta),
 \qquad
 \mathbf C(\theta)=\mathbf Q(\theta)^\trans
                     \mathbf C(0)\mathbf Q(\theta).
 \label{eq:basisrotation}
\end{equation}
The deformation gradient has one reference-basis index and therefore contains
one factor of \(\mathbf Q\).  The metric has two reference-basis indices and
contains two factors.  Its principal axes are directors rather than oriented
vectors, so its component orbit closes after a basis rotation of \(\pi\), while
the component matrix \(\mathbf F\) closes only after \(2\pi\).

For the example proposed in the question,
\begin{equation}
 \mathbf F(0)=\begin{pmatrix}1&0.4\\0&1\end{pmatrix},
 \qquad
 \mathbf C(0)=\begin{pmatrix}1&0.4\\0.4&1.16\end{pmatrix},
\end{equation}
one obtains
\begin{align}
 \mathbf F(\pi/2)&=
 \begin{pmatrix}0.4&-1\\1&0\end{pmatrix},
 &
 \mathbf C(\pi/2)&=
 \begin{pmatrix}1.16&-0.4\\-0.4&1\end{pmatrix},
 \\
 \mathbf F(\pi)&=-\mathbf F(0),
 &
 \mathbf C(\pi)&=\mathbf C(0).
\end{align}
Figure~\ref{fig:basisrotation} shows the corresponding basis vectors and the
closed orbit of \(\mathbf C\) in the Poincare-disk coordinates used by the
visualization code.  At \(\theta=\pi/2\), \(\mathbf C\) has moved to the
opposite elastic quadrant.  At \(\theta=\pi\), the oriented basis vectors have
reversed, but the metric has returned to its initial value.

\begin{figure}[H]
 \centering
 \includegraphics[width=0.82\textwidth]{figures/reference_basis_rotation.pdf}
 \caption{A passive rotation of the reference basis.  The top row shows the
 columns of \(\mathbf F(\theta)\).  The bottom panel shows that
 \(\mathbf C(\theta)\) completes a closed component-space orbit while the
 oriented columns of \(\mathbf F\) have completed only a half turn.  This is
 the same $90^\circ$ intermediate-basis relation that connects the short and
 long elastic-reduction decompositions.}
 \label{fig:basisrotation}
\end{figure}
\FloatBarrier

\section{Interpretation and conclusion}

Comparison with the original deformation gives two answers, depending on
whether the intermediate frame is part of the state.  If its axes are inherited
from the original reference frame and are required to remain nearby, the short
path is clearly preferable.  Its plastic factor has the same positive
horizontal-shear form as \(\mathbf F\), differs from \(\mathbf F\) by only
\(0.4\), and leaves an elastic factor only \(0.4\) from \(\mathbf I\).  The
long path inserts a quarter turn into the intermediate frame and is much less
compatible in these fixed components.  Equations~\eqref{eq:alignedplasticdistance}
and~\eqref{eq:fourclassscores} make this preference quantitative.

This suggests a useful selection rule: among admissible reductions, choose the
one that minimizes \(\lVert\mathbf F_E-\mathbf I\rVert_{\Frob}\), or,
equivalently for the present example, the one whose \(\mathbf F_P\) remains
closest to the total \(\mathbf F\) in the anchored frame.  In an incremental
calculation, the more robust version is to choose the representative closest
to the previously selected intermediate orientation.  This suppresses
unnecessary \(90^\circ\) or \(180^\circ\) jumps and makes the short path a
natural continuity choice.

If the intermediate axes are not anchored, the horizontal-versus-vertical
distinction is only a change of components.  Applying the admissible rotation
\(\mathbf Q\) turns the long-path \(\mathbf F_P\) exactly into the short-path
\(\mathbf F_P\).  Their objective elastic and plastic metrics, singular
values, energy, and Cauchy stress remain identical.  Consequently the original
\(\mathbf F\) alone cannot create an objective preference between the two
quadrants; it can select a preferred \emph{representative} only after an
orientation or continuity convention has been supplied.

A constitutive model can turn that convention into a physical preference by
carrying plastic orientation or labeled slip systems as history variables.
Alternatively, if every unit-shear event has the same positive dissipation,
the BFS word length is itself a cost and independently favors the two-step
short route over the three-step long route.  Thus the short quadrant is
defensible both because it preserves the anchored intermediate orientation and
because it has lower path cost.

\end{document}
